\section{High-Performance Implementation Details}
\label{sec:implementation}

We construct our routing engine in C++ using data-oriented design to maximize hardware efficiency~\cite{drepper2007memory}. We pack the road network topology into a Compressed Sparse Row (CSR) layout. To model turn restrictions without conditional jumps in the hot loop, we apply dual graph line inversion where road segments become nodes and turns become edges, mapping turn penalties to static edge weights. We apply a topological node reordering based on the Morton spatial curve (Z-curve) during graph compilation to assign adjacent addresses to spatial neighbors, letting the prefetcher load adjacent nodes into L1/L2 caches before access. 

We represent the dynamic traffic states using an Array of Structures (AoS) format. 
We align each segment state structure to 64-byte boundaries using the \verb|alignas(64)| directive. 
This alignment prevents cache line splitting and guarantees that a single cache transaction retrieves the node state. 
To eliminate runtime weight calculation overhead, we pre-compute and store link travel times directly within the graph memory buffers. The routing loop reads these integer edge weights as static values, avoiding floating-point division and BPR function evaluation in the hot path.

We implement a cache-conscious SIMD 8-ary priority heap to reduce tree traversal depth. 
The 8-ary structure matches the width of the AVX2 vector register: eight 32-bit child keys occupy exactly 256 bits, saturating one SIMD load instruction with no wasted lanes. 
We vectorize the child node search step using AVX2 intrinsics, aligning with hardware-level optimization guidelines for SIMD registers, data packing, and alignment on x86 architectures~\cite{kusswurm2014modern}. 
The child search routine uses the \verb|_mm256_min_epu32| instruction to locate the minimum among all eight child keys within a single 256-bit register, representing an instruction-level parallel operation executing on one CPU core. 
This instruction eliminates scalar branch instructions during the down-heap traversal and reduces the hot-loop critical path. This design explicitly targets modern x86-64 desktop and server architectures supporting AVX2 instructions, trading platform portability (e.g., to ARM or Apple Silicon architectures) for maximum low-level vectorization throughput.
 

We apply joint word encoding to represent heap elements. 
We pack the 29-bit distance key and the 3-bit local child index into a single 32-bit word. 
The horizontal SIMD minimum then returns the packed winner, so the minimum child key and its position decode in two bitwise operations with zero additional comparisons. 
This format also avoids pointer indirection. 

We precompute the ALT landmark distance table with exactly sixteen landmarks. 
Each landmark stores one 32-bit distance value; sixteen entries occupy $16 \times 4 = 64$ bytes, which fits exactly one CPU cache line. 
The ALT heuristic evaluation loop loads all sixteen landmark bounds in a single cache transaction and processes them without L1 misses. 

We utilize Target Register Caching to optimize coordinate access during A* and ALT searches. 
We load destination coordinates and landmark bounds into CPU registers before entering the main loop. 
This step reduces register pressure and eliminates stack spill operations. 

Metropolitan graphs containing 1.2 million nodes require rapid state cleanup between queries. 
Re-initializing the distance tracking array requires $O(V)$ operations, which dominates execution time for localized queries. 
We implement an $O(\text{active})$ partial fast-clearing mechanism. 
We maintain a thread-local stack of visited node indices during traversal. 
After a query terminates, we iterate over the active stack to reset the modified distance values to infinity. 
This approach restricts cleanup costs to the size of the relaxed search space.


\subsection{Precomputation and Landmarked Pre-Baked Data}
To support fast ALT queries, we perform offline precomputation of landmark distances. Our preprocessing pipeline implements a three-level data flow. First, raw road geometries are compiled into an Edge-Based Graph representation, modeling turn penalties statically as edge weights. Second, the landmark selection routine generates 128 candidate landmarks using MinMax (Farthest) and Avoid heuristics. We then select exactly sixteen landmarks from this candidate set using the MaxCover stochastic optimization algorithm~\cite{fuchs2010preprocessing}.
Third, we compute the shortest-path distance vectors from all landmarks to every node using Dijkstra searches. These precomputed distance matrices are stored as flat contiguous tables in memory. At query time, the ALT heuristic evaluation reads all sixteen precomputed 32-bit bounds in a single 64-byte L1 cache transaction, guaranteeing zero cache miss penalty in the routing hot loop.
